\documentclass[../main.tex]{subfiles}
\begin{document}
\section*{第1問}

% (原稿において第1問の解答は空欄)

\newpage
\section*{第2問}

{\bf [解]}

$P$から$\alpha$に下ろした垂足を$H$とし，$|PH|=h$とおく．また球の中心を$O$とする．

\begin{figure}[htb]
  \centering
  \begin{tikzpicture}[scale=2.2, >=stealth, every node/.style={font=\small}]

    % --- 1. 定数・幾何パラメータの定義 ---
    % 球の中心 O を原点とする (半径 R = 1.0)
    % 接点の高さ: z = cos\theta = 0.6
    % 接円の半径: r_c = sin\theta = 0.8
    % 頂点 P の高さ: z = -1 / cos\theta = -1 / 0.6 ≈ -1.667
    % 視線角に応じた接円楕円の短軸スケール比: b = 0.22
    \def\R{1.0}
    \def\zH{-0.6}       % スケッチに合わせて P を下向き (-z 方向) に配置
    \def\zP{-1.75}      % 頂点 P
    \def\rH{0.8}        % 平面 \alpha における接円の半径
    \def\bH{0.20}       % 楕円の短軸半径

    % 座標定義
    \coordinate (O) at (0, 0);          % 球の中心
    \coordinate (H) at (0, \zH);        % 垂足（接円の中心）
    \coordinate (P) at (0, \zP);        % 円錐の頂点
    \coordinate (Q_L) at (-\rH, \zH);   % 接円の左端
    \coordinate (Q_R) at (\rH, \zH);    % 接円の右端

    % --- 2. 円錐の母線（延長部を含む） ---
    % P から接点 Q_L, Q_R を通って上方に広がる母線
    \draw[thick, black] (P) -- ($ (P)!1.55!(Q_L) $);
    \draw[thick, black] (P) -- ($ (P)!1.55!(Q_R) $);

    % --- 3. 球の輪郭線 ---
    \draw[thick, black] (O) circle (\R);

    % --- 4. 平面 \alpha の接円（切り口の楕円） ---
    % 奥側 (破線)
    \draw[thick, dashed, teal!85!black]
    (Q_R) arc (0:180:{\rH} and {\bH});

    % 手前側 (実線)
    \draw[thick, teal!85!black]
    (Q_L) arc (180:360:{\rH} and {\bH});

    % --- 5. 中心軸・垂線 (O - H - P) ---
    % 球内部・奥側を通るため破線で描画
    \draw[dashed, gray!80!black, thick] (O) -- (P);

    % --- 6. 主要点のプロット ---
    \fill (O) circle (0.7pt);
    \fill (H) circle (0.7pt);
    \fill (P) circle (0.8pt) node[below=2pt] {$P$};

    % --- 7. ラベル・引き出し線 ---
    % 球の中心 O
    \draw[<-, thin, gray!80!black] (0.02, 0.05) -- (0.5, 0.6)
    node[above right, font=\footnotesize] {$O$};

    % 垂足 H
    \node[above right=1pt, font=\footnotesize] at (H) {$H$};

    % 平面 \alpha (接円)
    \draw[<-, thin, teal!85!black] (0.45, \zH - 0.12) -- (0.85, \zH - 0.18)
    node[right, font=\footnotesize, teal!90!black] {$\alpha$};

  \end{tikzpicture}
  \caption{円錐に内接する球と接円（平面 $\alpha$），および対称軸 $OP$ の3次元配置}
  \label{fig:sphere_cone_3d}
\end{figure}

この立体を$O$，$P$を含む平面で切ると，切断面は対称性から下図のようになる．

\begin{figure}[htb]
  \centering
  \begin{tikzpicture}[scale=2.8, >=stealth, every node/.style={font=\small}]

    % --- 1. 定数・幾何パラメータ定義 ---
    % 説明図として見やすい角度: cos\theta = 0.6, sin\theta = 0.8 (\theta ≈ 53.13°)
    \def\cVal{0.6}
    \def\sVal{0.8}
    \pgfmathsetmacro{\xP}{1.0 / \cVal} % 1 / cos\theta ≈ 1.6667
    \def\thetaDeg{53.1301}

    % 主要点の座標定義
    \coordinate (O) at (0, 0);
    \coordinate (P) at (\xP, 0);
    \coordinate (Q) at (\cVal, \sVal);
    \coordinate (Qprime) at (\cVal, -\sVal);
    \coordinate (H) at (\cVal, 0);

    % --- 2. 立体 D の断面（斜線ハッチング領域） ---
    % 円錐の母線 PQ, PQ' と 球欠の弧 Q' -> (1,0) -> Q で囲まれた領域
    \fill[pattern=north east lines, pattern color=blue!60, opacity=0.7]
    (Q) -- (P) -- (Qprime) arc (-\thetaDeg:\thetaDeg:1) -- cycle;

    % --- 3. 単位円（球の断面）の描画 ---
    \draw[thick, gray!70] (O) circle (1);

    % --- 4. 座標軸 ---
    \draw[->, thick] (-1.25, 0) -- (2.05, 0) node[right] {$x$};
    \draw[->, thick] (0, -1.25) -- (0, 1.35) node[above] {$y$};
    \node[below left=2pt] at (O) {$O$};

    % --- 5. 接線・母線および平面 \alpha (線分 QQ') ---
    % 接線（円錐の母線）
    \draw[very thick, red!85!black] ($ (P)!1.15!(Q) $) -- (P) -- ($ (P)!1.15!(Qprime) $);

    % 平面 \alpha の断面 (直線 x = \cos\theta)
    \draw[thick, dashed, teal!85!black] (\cVal, -1.15) -- (\cVal, 1.15)
    node[above, font=\footnotesize] {平面 $\alpha: x = \cos\theta$};

    % 半径 OQ
    \draw[thick, gray!80!black] (O) -- (Q);

    % --- 6. 直角マーク ---
    % 接点 Q における接線と半径の垂直マーク (OQ ⊥ QP)
    % 方向: OQ 方向 (\cVal, \sVal), QP 方向 (\cVal - \xP, \sVal)
    \draw[gray!80!black]
    ($ (Q) + (-0.06*\cVal, -0.06*\sVal) $) --
    ($ (Q) + (-0.06*\cVal - 0.08*\sVal, -0.06*\sVal + 0.08*\cVal) $) --
    ($ (Q) + (-0.08*\sVal, 0.08*\cVal) $);

    % 垂足 H における垂直マーク
    \draw[gray!80!black]
    ($ (H) + (0.06, 0) $) -- ($ (H) + (0.06, 0.06) $) -- ($ (H) + (0, 0.06) $);

    % --- 7. 角度 \theta ---
    \draw[purple!80!black, thick] (0.28, 0) arc (0:\thetaDeg:0.28);
    \node[purple!90!black, font=\footnotesize] at (0.42, 0.16) {$\theta$};

    % --- 8. 主要点のプロットとラベル ---
    \fill (O) circle (0.6pt);
    \fill (P) circle (0.7pt) node[below right=1pt, red!85!black] {$P\Bigl(\frac{1}{\cos\theta}, 0\Bigr)$};
    \fill (Q) circle (0.7pt) node[above right=1pt] {$Q(\cos\theta, \sin\theta)$};
    \fill (Qprime) circle (0.7pt) node[below right=1pt] {$Q'$};
    \fill (H) circle (0.7pt) node[below left=1pt] {$H(\cos\theta, 0)$};
    \fill (1, 0) circle (0.6pt) node[below right=1pt, font=\scriptsize] {$1$};

    % --- 9. 長さの寸法線 (高さ h = |PH|) ---
    %\draw[<->, blue!85!black, thick] ($ (H) + (0, -0.22) $) -- ($ (P) + (0, -0.22) $)
    %  node[midway, fill=white, inner sep=1pt, font=\scriptsize] {$h = \frac{1}{\cos\theta} - \cos\theta$};
    \draw[dotted, gray!60] (H) -- ($ (H) + (0, -0.25) $);
    \draw[dotted, gray!60] (P) -- ($ (P) + (0, -0.25) $);

    % 半径 1 の注記
    \node[gray!70!black, font=\scriptsize, above left] at ($ (O)!0.5!(Q) $) {$1$};

  \end{tikzpicture}
  \caption{球と円錐の対称軸を含む断面，および回転して立体 $D$ を得る領域（斜線部）}
  \label{fig:sphere_cone_section}
\end{figure}

ここで$O$を原点，$OP$を$x$軸正方向とする座標平面を定め，接点を$Q$として$\angle POQ=\theta$とする．ただし$0<\theta<\dfrac{\pi}{2}$とする．

$P$は円の接線$\cos\theta x+\sin\theta y=1$と$x$軸の交点だから，$P(1/\cos\theta,0)$となる．従って
\begin{align}
  |PH| = h = |OP| - |OH| =\dfrac{1}{\cos\theta}-\cos\theta \label{eq:1}
\end{align}
である．$D$は図の斜線部分を$x$軸のまわりに回転した立体であるから，その体積$V_D$は円錐から球欠を取り除いたもので
\begin{align}
  V_D
  &=\dfrac{1}{3}\pi \sin^2\theta h-\pi\int_{\cos\theta}^{1}(1-x^2)\,dx\\
  &=\dfrac{1}{3}\pi \sin^2\theta \Bigl(\dfrac{1}{\cos\theta}-\cos\theta\Bigr)
  -\pi\Bigl[\dfrac{2}{3}-\cos\theta+\dfrac{1}{3}\cos^3\theta\Bigr]\\
  &=\dfrac{\pi}{3}\Bigl[\dfrac{\sin^2\theta(1-\cos^2\theta)}{\cos\theta}+3\cos\theta-\cos^3\theta-2\Bigr]\\
  &=\dfrac{\pi}{3}\dfrac{\cos^2\theta-2\cos\theta+1}{\cos\theta} \label{eq:2}
\end{align}
となる．
これが球の体積の半分$\dfrac{2\pi}{3}$に等しいので，
\begin{align}
  V_D=\dfrac{2}{3}\pi
  \quad\Longleftrightarrow\quad
  \dfrac{\cos^2\theta-2\cos\theta+1}{\cos\theta}=2
  \quad\Longleftrightarrow\quad
  \cos\theta=2-\sqrt{3}\qquad(0<\cos\theta<1) \label{eq:3}
\end{align}
となる．このとき$K$の体積$V_K$は
\begin{align}
  V_K=\pi\int_{\cos\theta}^{1}(1-x^2)\,dx
  =\dfrac{2}{3}\pi\Bigl(11-6\sqrt{3}\Bigr)
\end{align}
である．

\newpage
\section*{第3問}

{\bf [解]}

漸化式
\begin{align}
  \begin{cases}
    u_1 = 2 \\
    u_2 = a^2+2 \\
    u_{n+2} = au_{n}-u_{n+1}
  \end{cases} \label{eq:1}
\end{align}
について考える．
初めの4項は
\begin{align}
  u_1=2，\quad u_2=a^2+2，\quad u_3=2a-a^2-2，\quad
  u_4=a^3+a^2+2 \label{eq:2}
\end{align}
となる．以下，合同式はすべて法4で考える
\cref{eq:2}において$a\equiv1$のとき$u_4\equiv0$となるので不適である．
従って$a\equiv0,2,-1$が必要である．以下各ケースについて個別に考える．

\subparagraph{$a\equiv0$のとき．}

すべての$n$について$u_n\not\equiv0$となることを数学的帰納法で示す．
$n=1,2$では$u_1\equiv u_2\equiv2$である．
そこで$k\in\mathbb{N}$に対して$n=k,k+1$での成立を仮定すれば，漸化式\cref{eq:1}から
\begin{align}
  u_{k+2}\equiv au_k - u_{k+1} \equiv -u_{k+1}\not\equiv0
\end{align}
となり，$n=k+2$でも成立する．従って数学的帰納法によりすべての$n$で$u_n\not\equiv0$であり，$4$の倍数でない．

\subparagraph{$a\equiv2$のとき．}
すべての$n$について$u_n\equiv2$となることを数学的帰納法で示す．
$n=1,2$では\cref{eq:2}より成立する．
そこで$k\in\mathbb{N}$に対して$n=k,k+1$での成立を仮定すれば，漸化式\cref{eq:1}から
\begin{align}
  u_{k+2} \equiv au_k - u_{k+1} \equiv 2u_k-u_{k+1}\equiv 4-2 \equiv 2
\end{align}
となり，$n=k+2$でも成立する．従って数学的帰納法によりすべての$n$で$u_n\equiv2$である．よって全ての$u_n$は$4$の倍数でない．

\subparagraph{$a\equiv-1$のとき．}
$u_n$は法4の剰余類で$(2,-1,-1)$の周期数列となることを数学的帰納法で示す．
$n=1,2,3$では実際に$(u_1,u_2,u_3)\equiv(2,-1,-1)$である．
ある$k\geqq1$について，$u_{3k-2},u_{3k-1},u_{3k}$での成立を仮定すると，漸化式から
\begin{align}
  u_{3k+1}&\equiv au_{3k-1} - u_{3k}   \equiv-(-1)-(-1)\equiv2，\\
  u_{3k+2}&\equiv au_{3k}   - u_{3k+1} \equiv -2-(-1)\equiv-1，\\
  u_{3k+3}&\equiv au_{3k+1} - u_{3k+2} \equiv -(-1)-2\equiv-1
\end{align}
となる．従って数学的帰納法により$(2,-1,-1)$の周期数列となり，$4$の倍数となる項は存在しない．
\casesend

以上より，必要十分条件は
\begin{align}
  a\equiv-1,0,2\pmod{4}
\end{align}
である．

\newpage
\section*{第4問}

{\bf [解]}

$O$を原点とする座標平面で考える．対称性から$P(1,0)$としてよい．
$O$を原点とする半径$1$の円を$C_1$，$P$を中心とする半径$2/\sqrt{3}$の円を$C_2$とすると，これらの方程式は
\begin{align}
  \begin{cases}
    C_1:x^2+y^2=1
    C_2:(x-1)^2+y^2=\dfrac{4}{3}
  \end{cases} \label{eq:1}
\end{align}
であり，$Q$，$R$は$C_2$上にある．

また，$C_1$，$C_2$の二つの交点$A,B$とすると，この$x$座標は$1/3$である．$A$を$y$座標が正の方とすると
\begin{align}
  A\Bigl(\dfrac{1}{3},\dfrac{2\sqrt{2}}{3}\Bigr) \label{eq:2}
\end{align}
となる．$\angle APO=\alpha$（$0<\alpha<\pi/2$）とすると
\begin{align}
  \tan\alpha=\dfrac{2\sqrt{2}/3}{2/3}=\sqrt{2} \label{eq:3}
\end{align}
である．

$\angle OPQ=\theta$とおくと，$Q$，$R$の座標は
\begin{align}
  \begin{cases}
    Q&=\Bigl(1-\dfrac{2}{\sqrt{3}}\cos\theta，\dfrac{2}{\sqrt{3}}\sin\theta\Bigr)，\\
    R&=\Bigl(1-\dfrac{2}{\sqrt{3}}\cos(\theta-\pi/3)，
    \dfrac{2}{\sqrt{3}}\sin(\theta-\pi/3)\Bigr)
  \end{cases} \label{eq:4}
\end{align}
となる．これらの様子を下図に示す．

\begin{figure}[htb]
  \centering
  \begin{tikzpicture}[scale=2.8, >=stealth, every node/.style={font=\small}]

    % --- 1. 定数・幾何パラメータの定義 ---
    % C1: 原点 O(0,0), 半径 R1 = 1
    % C2: 中心 P(1,0), 半径 R2 = 2/sqrt(3) ≈ 1.1547
    % 交点 A, B: x = 1/3, y = ±2*sqrt(2)/3 ≈ ±0.9428
    % \cos\alpha = 1/\sqrt{3} => \alpha ≈ 54.7356°
    \def\Rone{1.0}
    \pgfmathsetmacro{\Rtwo}{2.0 / sqrt(3)}       % 2/sqrt(3) ≈ 1.1547
    \def\alphaDeg{54.7356}                       % \alpha
    \pgfmathsetmacro{\angA}{180.0 - \alphaDeg}   % A の偏角 ≈ 125.2644°
    \pgfmathsetmacro{\angB}{180.0 + \alphaDeg}   % B の偏角 ≈ 234.7356°

    % 一般の \theta (例: \theta = 46° と設定，許容範囲 -\alpha+60° <= \theta <= \alpha を満たす)
    \def\thetaDeg{46.0}
    \pgfmathsetmacro{\angQ}{180.0 - \thetaDeg}                 % Q の偏角 (P基準): 134°
    \pgfmathsetmacro{\angR}{180.0 - (\thetaDeg - 60.0)}        % R の偏角 (P基準): 194°

    % 主要座標の定義
    \coordinate (O) at (0, 0);
    \coordinate (P) at (1, 0);
    \coordinate (A) at (1.0/3.0, {2.0*sqrt(2)/3.0});
    \coordinate (B) at (1.0/3.0, {-2.0*sqrt(2)/3.0});

    % 一般の Q, R の座標
    \coordinate (Q) at ($ (P) + (\angQ:\Rtwo) $);
    \coordinate (R) at ($ (P) + (\angR:\Rtwo) $);

    % --- 2. 単位円 C_1 内部に存在する許容弧 AB の強調 ---
    \draw[line width=1.8pt, teal!80!black]
    ([shift=(\angA:\Rtwo)]P) arc (\angA:\angB:\Rtwo);

    % --- 3. 円 C_1, C_2 の全体描画 ---
    % 単位円 C_1
    \draw[thick, gray!75!black] (O) circle (\Rone);
    \node[gray!75!black, font=\footnotesize, above left] at (-0.75, 0.75) {$C_1: x^2+y^2=1$};

    % 円 C_2 (C_1 外部は薄い破線で表示)
    \draw[thick, dashed, teal!50]
    ([shift=(\angB:\Rtwo)]P) arc (\angB:{360+\angA}:\Rtwo);
    \node[teal!85!black, font=\footnotesize, right] at ($ (P) + (35:\Rtwo) $) {$C_2: (x-1)^2+y^2=\frac{4}{3}$};

    % --- 4. 座標軸 ---
    \draw[->, thick] (-1.35, 0) -- (2.45, 0) node[right] {$x$};
    \draw[->, thick] (0, -1.35) -- (0, 1.35) node[above] {$y$};
    \node[below left=2pt] at (O) {$O$};

    % --- 5. 正三角形 PQR（一般の \theta のとき）---
    \fill[blue!15, opacity=0.7] (P) -- (Q) -- (R) -- cycle;
    \draw[very thick, blue!85!black] (P) -- (Q) -- (R) -- cycle;

    % 線分 OQ, OR
    \draw[thick, densely dashed, blue!75!black] (O) -- (Q);
    \draw[thick, densely dashed, blue!75!black] (O) -- (R);

    % 正三角形の頂角 (60°)
    \draw[thin, blue!90!black] ($ (P) + (\angQ:0.35) $) arc (\angQ:\angR:0.35);
    \node[blue!90!black, font=\tiny, left=1pt] at ($ (P) + ({0.5*(\angQ+\angR)}:0.42) $) {$\frac{\pi}{3}$};

    % --- 6. 角度 \alpha, \theta の表示 ---
    % 線分 PA, PB
    \draw[thick, gray!70] (P) -- (A);
    \draw[thick, gray!70] (P) -- (B);

    % \angle APO = \alpha (P を中心とする負の x 軸からの限界角)
    \draw[thin, purple!80!black] ($ (P) + (180:0.52) $) arc (180:\angA:0.52);
    \node[purple!90!black, font=\tiny] at ($ (P) + (155:0.60) $) {$\alpha$};

    % \theta = \angle OPQ (一般角)
    \draw[thin, orange!90!black] ($ (P) + (180:0.68) $) arc (180:\angQ:0.68);
    \node[orange!90!black, font=\tiny] at ($ (P) + ({180 - 0.5*\thetaDeg}:0.78) $) {$\theta$};

    % --- 7. 交点 A, B の座標射影補助線 ---
    \draw[dotted, gray!70] (A) -- (1.0/3.0, 0) -- (B);
    \draw[dotted, gray!70] (A) -- (0, {2.0*sqrt(2)/3.0});
    \fill (1.0/3.0, 0) circle (0.5pt);
    \node[below=2pt, font=\scriptsize] at (1.0/3.0, 0) {$\frac{1}{3}$};
    \node[left=2pt, font=\tiny] at (0, {2.0*sqrt(2)/3.0}) {$\frac{2\sqrt{2}}{3}$};

    % --- 8. 主要点のプロットとラベル ---
    \fill (P) circle (0.7pt) node[below right=1pt] {$P(1,0)$};
    \fill (A) circle (0.7pt) node[above=2pt] {$A\Bigl(\frac{1}{3}, \frac{2\sqrt{2}}{3}\Bigr)$};
    \fill (B) circle (0.7pt) node[below=2pt] {$B$};
    \fill[blue!85!black] (Q) circle (0.7pt) node[above left=1pt] {$Q$};
    \fill[blue!85!black] (R) circle (0.7pt) node[below left=1pt] {$R$};
    \fill (O) circle (0.6pt);

  \end{tikzpicture}
  \caption{円 $C_1, C_2$ の許容弧 $AB$ と一般の角 $\theta$ における正三角形 $PQR$}
  \label{fig:circle_C1_C2_general_theta}
\end{figure}

さて，$Q,R$が単位円の内部または周上にある条件は
\begin{align}
  -\alpha\leqq\theta-\pi/3<\theta\leqq\alpha
  \quad\Longleftrightarrow\quad
  -\alpha+\pi/3\leqq\theta\leqq\alpha \label{eq:5}
\end{align}
である．以下簡単のために$c=\cos\theta$，$s=\sin\theta$とすると$|OQ|^2,|OR|^2$は
\begin{align}
  |OQ|^2
  &=\Bigl(1-\dfrac{2}{\sqrt{3}}c\Bigr)^2
  +\Bigl(\dfrac{2}{\sqrt{3}}s\Bigr)^2
  =\dfrac{7}{3}-\dfrac{4}{\sqrt{3}}c，\label{eq:6}\\
  |OR|^2&=\dfrac{7}{3}-\dfrac{4}{\sqrt{3}}\cos(\theta-\pi/3)
  =\dfrac{7}{3}-\dfrac{4}{\sqrt{3}}\Bigl(\dfrac{1}{2}c+\dfrac{\sqrt{3}}{2}s\Bigr)．\label{eq:7}
\end{align}
である．従って$f(\theta)=|OQ|^2+|OR|^2$として\cref{eq:6,eq:7}を代入すると
\begin{align}
  \begin{split}
    f(\theta)&=\dfrac{14}{3} - \dfrac{4\sqrt{3}}{3}\Bigl(\dfrac{3}{2}c+\dfrac{\sqrt{3}}{2}s\Bigr)\\
    &=\dfrac{14}{3}-4\sin(\theta+\pi/3)
  \end{split} \label{eq:8}
\end{align}
である．\cref{eq:5}より
\begin{align}
  -\alpha+2\pi/3\leqq\theta+\pi/3\leqq\alpha+\pi/3
\end{align}
であり，この区間は$\theta+\dfrac{\pi}{3}=\dfrac{\pi}{2}$を含む．
よってこの時$\theta=\dfrac{\pi}{6}$で最小となり，
最大値は区間の両端の値のうち大きい方である．
最小値は
\begin{align}
  \min f=f\left(\dfrac{\pi}{6}\right) = \dfrac{14}{3}-4=\dfrac{2}{3}
\end{align}
である．また\cref{eq:3}から$\sin\alpha=\sqrt{2/3}$，$\cos\alpha=1/\sqrt{3}$なので
\begin{align}
  f(\alpha)&=\dfrac{14}{3}
  -4\Bigl(\dfrac{1}{2}\sqrt{\dfrac{2}{3}}+\dfrac{\sqrt{3}}{2}\dfrac{1}{\sqrt{3}}\Bigr)
  =\dfrac{2}{3}(4-\sqrt{6})，\\
  f(\pi/3-\alpha)&=\dfrac{14}{3}-4\sin(2\pi/3-\alpha)
  =\dfrac{2}{3}(4-\sqrt{6})
\end{align}
となり両端での値は同じで，これらが同時に最大値となる．

以上から求める最大値と最小値は
\begin{align}
  \text{最大値 }\dfrac{2}{3}(4-\sqrt{6})，\qquad\text{最小値 }\dfrac{2}{3}
\end{align}
である．

\newpage
\section*{第5問}

{\bf [解]}

$f(x)$は3次式であり，条件(B)から，$a\neq0$として
\begin{align}
  f''(x)=2ax+2 \label{eq:1}
\end{align}
と書ける．積分して条件$f'(0)=0$を用いると
\begin{align}
  f'(x)=ax^2+2x \label{eq:2}
\end{align}
だから，さらに積分して実数$b$を用いて
\begin{align}
  f(x)=\dfrac{a}{3}x^3+x^2+b \label{eq:3}
\end{align}
と書ける．

一方で条件(A)より$f(x)-x$は$x=-t$，$x=0$を根にもつので
\begin{align}
  &
  \begin{cases}
    f(-t)+t = 0 \\
    f(0) = 0
  \end{cases} \\
  \therefore
  &
  \begin{cases}
    -\dfrac{a}{3}t^3+t^2+t+b=0 \\
    b=0
  \end{cases} \\
  \therefore
  &
  \begin{cases}
    t\left(-\dfrac{a}{3}t^2+t+1\right) = 0 \\
    b=0
  \end{cases} \label{eq:4}
\end{align}
である．$t>0$から第一式は
\begin{align}
  a=\dfrac{3(t+1)}{t^2} \label{eq:5}
\end{align}
となる．以上\cref{eq:4,eq:5}より$f(x)$は
\begin{align}
  f(x) = \dfrac{a}{3}x^3+x^2 \label{eq:6}
\end{align}
であり，$f(x)-x$を$g(x)$とおくと
\begin{align}
  g(x)
  &=f(x)-x \\
  &=\dfrac{t+1}{t^2}x^3+x^2-x\\
  &=\dfrac{1}{t^2}x(x+t)\Bigl((t+1)x-t\Bigr)
\end{align}
と因数分解できる．これを図示すると下図であり，求める面積$S(t)$は斜線部の領域の面積である．

\begin{figure}[htb]
  \centering
  \begin{tikzpicture}[scale=3.0, >=stealth, every node/.style={font=\small}]

    % --- 1. 定数・代表点のパラメータ定義 ---
    % t = 1 のとき:
    % a = 3*(1+1)/1^2 = 6  =>  f(x) = 2x^3 + x^2
    % d = t/(t+1) = 1/2 = 0.5
    % 交点: x = -1 (-t), x = 0 (原点), x = 0.5 (d)
    \def\tVal{1.0}
    \def\dVal{0.5}

    % --- 2. 求める面積 S(t) のハッチング領域 (0 <= x <= d) ---
    % 直線 y = x が上、曲線 y = f(x) が下
    \fill[pattern=north east lines, pattern color=blue!60, opacity=0.7]
    (0, 0) -- (\dVal, \dVal)
    -- plot[domain=\dVal:0, samples=60] (\x, {2.0*\x*\x*\x + \x*\x})
    -- cycle;

    % --- 3. 座標軸 ---
    \draw[->, thick] (-1.35, 0) -- (0.95, 0) node[right] {$x$};
    \draw[->, thick] (0, -1.35) -- (0, 1.15) node[above] {$y$};
    \node[below left=2pt] at (0, 0) {$O$};

    % --- 4. 直線 y = x ---
    \draw[thick, orange!90!black, domain=-1.25:0.85]
    plot (\x, {\x});
    \node[orange!90!black, font=\footnotesize, above right] at (0.75, 0.75) {$y = x$};

    % --- 5. 3次曲線 y = f(x) ---
    \draw[very thick, blue!85!black, domain=-1.18:0.65, samples=100]
    plot (\x, {2.0*\x*\x*\x + \x*\x});
    \node[blue!85!black, font=\footnotesize, left] at (0.62, {2.0*0.62^3 + 0.62^2}) {$y = f(x)$};

    % --- 6. 交点および射影補助線 ---
    % (i) 交点 P(-t, -t)
    \draw[dashed, gray!70] (-\tVal, 0) -- (-\tVal, -\tVal) -- (0, -\tVal);
    \fill[red!80!black] (-\tVal, -\tVal) circle (0.7pt);
    \node[above left=1pt, font=\footnotesize, red!80!black] at (-\tVal, -\tVal) {$P$};
    \node[above=2pt, font=\scriptsize] at (-\tVal, 0) {$-t$};
    \fill (-\tVal, 0) circle (0.5pt);

    % (ii) 原点 O(0, 0)
    \fill[black] (0, 0) circle (0.6pt);

    % (iii) 交点 Q(d, d)
    \draw[dashed, gray!70] (\dVal, 0) -- (\dVal, \dVal) -- (0, \dVal);
    \fill[red!80!black] (\dVal, \dVal) circle (0.7pt);
    \node[above left=1pt, font=\footnotesize, red!80!black] at (\dVal, \dVal) {$Q$};
    \node[below=2pt, font=\scriptsize] at (\dVal, 0) {$d = \frac{t}{t+1}$};
    \fill (\dVal, 0) circle (0.5pt);

  \end{tikzpicture}
  \caption{3次曲線 $y = f(x)$ と直線 $y = x$，および面積 $S(t)$ を与える領域（斜線部）}
  \label{fig:cubic_line_area_St}
\end{figure}

簡単のために$d=t/(t+1)$とおくと求める面積$S(t)$は
\begin{align}
  S(t)
  &=-\int_0^dg(x)\,dx\\
  &=\Bigl[-\dfrac{t+1}{4t^2}x^4-\dfrac{1}{3}x^3+\dfrac{1}{2}x^2\Bigr]_0^d\\
  &=-\dfrac{1}{4}\Bigl(\dfrac{1}{t}+\dfrac{1}{t^2}\Bigr)
  \Bigl(\dfrac{1}{1+1/t}\Bigr)^4
  -\dfrac{1}{3}\Bigl(\dfrac{1}{1+1/t}\Bigr)^3
  +\dfrac{1}{2}\Bigl(\dfrac{1}{1+1/t}\Bigr)^2
\end{align}
となる．よって$t\to\infty$の極限をとって
\begin{align}
  \lim_{t\to\infty}S(t)=\dfrac{1}{2}-\dfrac{1}{3}=\dfrac{1}{6}
\end{align}
である．

\newpage
\section*{第6問}

{\bf [解]}

題意から三角形$P_nP_{n+1}P_{n+2}$はすべて辺の長さ$1,a,\sqrt{1+a^2}$の相似である．
$\angle P_1P_0P_2=\alpha$，$\angle P_0P_1P_2=\beta$とする．

\begin{figure}[htb]
  \centering
  \begin{tikzpicture}[scale=4.5, >=stealth, every node/.style={font=\small}]

    % --- 1. 定数・幾何パラメータの定義 ---
    % a の値を設定（1 : a の比率がわかりやすい値として a = 1.35）
    \def\aVal{1.35}

    % P_0, P_1, P_2 の定義
    \coordinate (P0) at (1, 0);
    \coordinate (P1) at (0, \aVal);
    \coordinate (P2) at (0, 0);

    % P3: P2 から P0P1 への垂足
    \coordinate (P3) at ($ (P0)!(P2)!(P1) $);
    % P4: P3 から P1P2 (y軸) への垂足
    \coordinate (P4) at ($ (P1)!(P3)!(P2) $);
    % P5: P4 から P2P3 への垂足
    \coordinate (P5) at ($ (P2)!(P4)!(P3) $);
    % P6: P5 から P3P4 への垂足
    \coordinate (P6) at ($ (P3)!(P5)!(P4) $);
    % P7: P6 から P4P5 への垂足（収束方向の参考用）
    \coordinate (P7) at ($ (P4)!(P6)!(P5) $);

    % --- 2. 相似三角形の着色（第1世代と第2世代の対比） ---
    % \triangle P_0 P_1 P_2
    \fill[blue!8] (P0) -- (P1) -- (P2) -- cycle;
    % \triangle P_4 P_5 P_6（180度逆向きに相似）
    \fill[orange!25] (P4) -- (P5) -- (P6) -- cycle;

    % --- 3. 骨格線・垂線の描画 ---
    % 初期の直角三角形 P0P1P2
    \draw[very thick, black] (P0) -- (P1) -- (P2) -- cycle;

    % 垂線の軌跡 (P2 -> P3 -> P4 -> P5 -> P6 -> P7)
    \draw[thick, blue!85!black] (P2) -- (P3) -- (P4) -- (P5) -- (P6) -- (P7);

    % \triangle P_4P_5P_6 の斜辺 P4P5 以外の辺（補助線）
    \draw[dashed, orange!90!black, thick] (P4) -- (P6);

    % --- 4. 直角マーク (P2, P3, P4, P5, P6) ---
    % P2: (0,0)
    \draw[thin, gray!70] (0.06, 0) -- (0.06, 0.06) -- (0, 0.06);

    % P3: P2-P3 ⊥ P0-P1
    \coordinate (u3) at ($ (P1) - (P0) $);
    \coordinate (v3) at ($ (P2) - (P3) $);
    \draw[thin, gray!70]
    ($ (P3) + 0.045*({-1/\aVal}, 1) $) --
    ($ (P3) + 0.045*({-1/\aVal}, 1) + 0.045*(1, {1/\aVal}) $) --
    ($ (P3) + 0.045*(1, {1/\aVal}) $);

    % P4: P3-P4 ⊥ y軸
    \draw[thin, gray!70] ($ (P4) + (0.05, 0) $) -- ($ (P4) + (0.05, -0.05) $) -- ($ (P4) + (0, -0.05) $);

    % P5: P4-P5 ⊥ P2-P3
    \draw[thin, gray!70]
    ($ (P5) + 0.035*(1, {1/\aVal}) $) --
    ($ (P5) + 0.035*(1, {1/\aVal}) + 0.035*({-1/\aVal}, 1) $) --
    ($ (P5) + 0.035*({-1/\aVal}, 1) $);

    % P6: P5-P6 ⊥ P3-P4
    \draw[thin, gray!70] ($ (P6) + (0, 0.035) $) -- ($ (P6) + (-0.035, 0.035) $) -- ($ (P6) + (-0.035, 0) $);

    % --- 5. 角度 \alpha, \beta の表示 ---
    % \alpha = \angle P1 P0 P2
    \draw[purple!85!black, thick] ($ (P0) + (180:0.18) $) arc (180:{180 - atan(\aVal)}:0.18);
    \node[purple!90!black, font=\footnotesize] at ($ (P0) + (160:0.25) $) {$\alpha$};

    % \beta = \angle P0 P1 P2
    \draw[teal!85!black, thick] ($ (P1) + (-90:0.18) $) arc (-90:{atan(1/\aVal) - 90}:0.18);
    \node[teal!90!black, font=\footnotesize] at ($ (P1) + (-70:0.26) $) {$\beta$};

    % --- 6. 辺長ラベル ---
    \node[below=2pt] at ($ (P2)!0.5!(P0) $) {$1$};
    \node[left=2pt] at ($ (P2)!0.5!(P1) $) {$a$};
    \node[above right=1pt] at ($ (P0)!0.5!(P1) $) {$\sqrt{a^2+1}$};

    % --- 7. 主要点のプロットとラベル ---
    \fill (P0) circle (0.6pt) node[below right=1pt] {$P_0(1,0)$};
    \fill (P1) circle (0.6pt) node[above left=1pt] {$P_1(0,a)$};
    \fill (P2) circle (0.6pt) node[below left=1pt] {$P_2(0,0)$};
    \fill (P3) circle (0.5pt) node[above right=1pt] {$P_3$};
    \fill (P4) circle (0.5pt) node[left=1pt] {$P_4$};
    \fill (P5) circle (0.5pt) node[below=2pt] {$P_5$};
    \fill[red!85!black] (P6) circle (0.6pt) node[above=2pt, red!85!black] {$P_6$};

  \end{tikzpicture}
  \caption{繰り返される垂線と相似な直角三角形の連鎖（$\triangle P_0P_1P_2$ と $\triangle P_4P_5P_6$）}
  \label{fig:perpendicular_triangles}
\end{figure}

$x$軸を実軸，$y$軸を虚軸とする複素数平面で考える．
簡単のために
\begin{align}
  e(\theta)=\cos\theta+i\sin\theta
\end{align}
とおく．
ベクトル$\overrightarrow{P_nP_{n+1}}$を表す複素数を$q_n$とする．
題意より
\begin{align}
  \begin{cases}
    q_{2k+2}&=\dfrac{1}{\sqrt{a^2+1}}e(\pi-\beta)q_{2k+1}，\\
    q_{2k+1}&=\dfrac{a}{\sqrt{a^2+1}}e(\pi-\alpha)q_{2k}
  \end{cases} \label{eq:1}
\end{align}
だから$n$の偶奇にかかわらず
\begin{align}
  q_{n+2}=\dfrac{a}{a^2+1}e(\pi-\alpha)e(\pi-\beta)q_n \label{eq:2}
\end{align}
である．簡単のためにこの係数を
\begin{align}
  A=\dfrac{a}{a^2+1}e(\pi-\alpha)e(\pi-\beta) \label{eq:A-definition}
\end{align}
とおくと，\cref{eq:2}を繰り返し用いて
\begin{align}
  q_{2k}=A^kq_0，\qquad q_{2k+1}=A^kq_1 \label{eq:3}
\end{align}
を得る．初期条件から
\begin{align}
  q_0=-1+ai，\qquad q_1=-ai． \label{eq:4}
\end{align}
だから\cref{eq:3}に代入すると
\begin{align}
  q_{2k}=A^k(-1+ai)，\qquad q_{2k+1}=A^k(-ai) \label{eq:5}
\end{align}

ついで$A$を具体的に計算する．$\alpha+\beta=\pi/2$に注意して
\begin{align}
  e(\pi-\alpha)e(\pi-\beta)
  =e(2\pi-\alpha-\beta)=e(3\pi/2)=-i
\end{align}
だから，\cref{eq:2}に代入して
\begin{align}
  A=-ai/(a^2+1) \label{eq:6}
\end{align}
である．

点$P_n$を表す複素数を$p_n$とすると，$n\geqq1$について
\begin{align}
  p_{2n}
  &=p_0+\sum_{k=0}^{n-1}(q_{2k}+q_{2k+1})\\
  &=1-\sum_{k=0}^{n-1}A^k \quad (\because \text{\cref{eq:5}}) \\
  &=1-\dfrac{1-A^n}{1-A} \\
  &=\dfrac{-A+A^n}{1-A} \label{eq:7}
\end{align}
\begin{align}
  p_{2n+1}
  &=p_{2n}+q_{2n}
  =\dfrac{-A+A^n}{1-A}+A^n(-1+ai) \label{eq:8}
\end{align}
がなりたつ．

\paragraph{(1)}
$p_6$は\cref{eq:7}に$n=3$を代入して
\begin{align}
  \begin{split}
    p_6&=\dfrac{-A+A^3}{1-A}=-(1+A)A\\
    &=\dfrac{a^2}{(a^2+1)^2}+\dfrac{a}{a^2+1}i
  \end{split}
\end{align}
となる．従って
\begin{align}
  P_6\Bigl(\dfrac{a^2}{(a^2+1)^2}，\dfrac{a}{a^2+1}\Bigr)
\end{align}
である．

\paragraph{(2)}
$|A|=\dfrac{a}{a^2+1}<1$であるから，$n\to\infty$では$A^n\to 0$である．
従って\cref{eq:7,eq:8}より偶数番目，奇数番目ともに
\begin{align}
  \begin{split}
    \lim_{n\to\infty}p_n
    &= \dfrac{-A}{1-A}=\dfrac{ai}{a^2+1+ai}\\
    &=\dfrac{a^2+a(a^2+1)i}{(a^2+1)^2+a^2}
  \end{split}
\end{align}
となる．従って求める点は
\begin{align}
  \Bigl(\dfrac{a^2}{a^4+3a^2+1}，
  \dfrac{a(a^2+1)}{a^4+3a^2+1}\Bigr)
\end{align}
である．

{\bf [別解]}

$\triangle P_0P_1P_2$と$\triangle P_4P_5P_6$は180度逆向きで相似比は
$\bigl(a/(a^2+1)\bigr)^2$である．従って
\begin{align}
  \overrightarrow{P_6P_{10}}
  =-\Bigl(\dfrac{a}{a^2+1}\Bigr)^2\overrightarrow{P_2P_6}．
\end{align}
同様に繰り返して，$k\geqq1$について
\begin{align}
  \overrightarrow{P_{4k-2}P_{4k+2}}
  =\Bigl\{-\Bigl(\dfrac{a}{a^2+1}\Bigr)^2\Bigr\}^{k-1}
  \overrightarrow{P_2P_6}
\end{align}
となる．ここで
\begin{align}
  \overrightarrow{P_2P_6}
  =\Bigl(\dfrac{a^2}{(a^2+1)^2}，\dfrac{a}{a^2+1}\Bigr)
\end{align}
であるから
\begin{align}
  \begin{split}
    \overrightarrow{P_2P_{4n+2}}
    &=\sum_{k=1}^n
    \Bigl\{-\Bigl(\dfrac{a}{a^2+1}\Bigr)^2\Bigr\}^{k-1}
    \overrightarrow{P_2P_6}\\
    &\longrightarrow\dfrac{1}{1+\bigl(a/(a^2+1)\bigr)^2}
    \overrightarrow{P_2P_6}\\
    &=\Bigl(\dfrac{a^2}{a^4+3a^2+1}，
    \dfrac{a(a^2+1)}{a^4+3a^2+1}\Bigr)
  \end{split}
\end{align}
となる．相似比から各辺$|P_nP_{n+1}|$は0に収束するので，
残りの添字の点列も同じ点に収束する．

\newpage
\end{document}
